% Samenvatting, max 3500 tekens, Nederlands.

Deze masterproef beschrijft het ontwerp, de implementatie en de veldvalidatie van een draagbaar, SatNOGS-compatibel grondstation voor de AetherSpace CubeSat op 433\,MHz. Het station bestaat uit een Raspberry Pi 3A+ met twee subsystemen.

De rotator gebruikt een 3D-geprinte ASA-structuur met geïntegreerde geprinte lagerbanen en stalen kogels, gereduceerde DC-motoren (516:1 intern, tot 1376:1 totaal) en een op maat ontworpen RP2040-motorcontroller-PCB. Een PID-regellus haalt sub-0,5\degree\ elevatie- en 0--3\degree\ azimutfout, gevalideerd tijdens live passages. Elevatiekalibratie gebeurt via een ADXL345-accelerometer; azimut wordt gekalibreerd via het telefoonkompas of door op te starten vooraf op het noorden gericht.

Het RF-front-end is een 6-laags Pi HAT met twee TI CC1200 transceivers; het UHF-pad (432\,MHz) is bestukt met een matching network afgeleid van de TI CC1200EM-referentie, terwijl het VHF-pad (144/145\,MHz) gerouteerd is maar niet bestukt.

Een Python-stack op de Pi verzorgt passagevoorspelling, Dopplercorrectie, per-satelliet CC1200-herconfiguratie en indiening van SiDS-telemetrie. Een HTTPS-dashboard biedt laptop-gestuurd veldbeheer, met optionele telefoontoegang. Het station spreekt standaard EasyComm III via hamlib.

Over drie veldsessies werden ongeveer 210 satellietpassages gevolgd met een Siretta Oscar 44 Yagi. De CC1200-ontvangstketen is gevalideerd voor lokale testframes in korteafstandstests tussen twee CC1200-transceivers op de grond, maar geen protocol-geldige frames werden uit baan gedecodeerd. De link-budgetanalyse wijst een mastgemonteerde LNA aan als de belangrijkste resterende gevoeligheidskloof. De v2 MOTOR\_RF\_HAT (gecombineerde controller- en RF-print) is ontwerp-volledig maar bij indiening nog niet gefabriceerd.



